\section{Orders, thresholds, and vanishing}\label{sec:preliminaries}
Throughout the proof, $X$, $L$, and $n$ satisfy the hypotheses of
Theorem~\ref{thm:main}. Write
\[
 P=K_X+(n+1)L,\qquad R_j=H^0(X,jL),\qquad N_j=\dim R_j.
\]
Ampleness gives
\begin{equation}\label{eq:hilbert-X}
 N_j=\frac{L^n}{n!}j^n+O(j^{n-1}),\qquad L^n\geq1.
\end{equation}
All centers are \emph{closed} centers. Thus a divisor centered along
a subvariety $W$ is an admissible test at every point of $W$, not just
at its generic point.

\subsection{Weighted orders}
We use resolution and principalization in characteristic zero with a
prescribed simple normal crossing (SNC) boundary; see
\cite[Theorems~1.0.1 and~2.4.1]{wlodarczyk2005}.
Models used for vanishing are projective. We may resolve above an
already chosen model and retain any previously extracted prime by its
strict transform.

For a prime divisor $F$ on a smooth birational model $\pi:Y\to X$,
the log discrepancy is
\[
 \A(F)=1+\ord_F(K_{Y/X}).
\]
We permit positive multiples of $\ord_F$ and scale discrepancy by the
same multiple. Orders of sections mean orders of their effective zero
divisors; orders of ideals mean the minimum of the orders of local
generators. A local frame pulled back from $X$ may be used in these
calculations.

We also use weighted monomial orders. Choose a smooth projective
birational model and local regular coordinates adapted to the
\emph{entire} relative canonical divisor: every component of that
divisor meeting the coordinate neighborhood is a coordinate
hyperplane. All coordinate divisors used below, together with that
support, must form a simple normal crossing divisor. Choose a subset
$E_1,\ldots,E_r$ of the coordinate divisors and give it positive real
weights $u_1,\ldots,u_r$; the remaining coordinates have weight zero.
The order of a germ is the least weight of a nonzero
monomial in their local equations, with coefficients in the remaining
coordinates, and
\begin{equation}\label{eq:monomial-discrepancy}
             \A(v)=\sum_{j=1}^r u_j\A(E_j).
\end{equation}
A coordinate divisor absent from the relative canonical divisor has
log discrepancy one. We use a chosen irreducible component of the
positive-support stratum, whose generic point avoids all zero-weight
coordinate divisors; its closed image is the center on $X$. This
adaptation condition is essential for
\eqref{eq:monomial-discrepancy}. The test with all weights zero is
excluded.

For positive rational weights this is a divisorial order up to scale.
Indeed, clearing denominators and extracting the corresponding
primitive ray is a toroidal modification in regular parameters. The
relative canonical order in those parameters is the sum of the
primitive weights minus one. Adding the pullback of the previous
relative canonical divisor gives
\eqref{eq:monomial-discrepancy}. The extraction can be retained on a
smooth projective model by resolution. Approximating weights on a
fixed support by positive rational weights approximates both
discrepancy and the orders of any fixed finite collection of sections.
For the latter assertion only finitely many coordinatewise minimal
exponents need be considered.

\begin{lemma}[Local computation of a monomial order]\label{lem:monomial-local}
A monomial order with fixed positive support can be computed in the
Taylor ring at any smooth point of its open stratum. Its value is
unchanged by replacing each positive-weight coordinate by a unit
multiple or by changing the complementary zero-weight coordinates.
\end{lemma}
\begin{proof}
For a positive threshold $c$, the corresponding ideal is generated by
the monomials in the positive-weight parameters of weight at least
$c$. These generators are preserved, up to units, under the stated
changes. In the completed local ring the ideal is extended from a
finite-colength monomial ideal in the positive-weight variables. It
is therefore primary to their intersection. Localization at the
generic point of that stratum followed by contraction does not change
membership. The same conclusion in the regular local ring follows
from faithful flatness of completion. In particular a coefficient
which vanishes at the chosen closed point, but not identically on the
stratum, is still a nonzero series in zero-weight variables; it does
not raise the old weighted order.
\end{proof}

\begin{lemma}[SNC retraction]\label{lem:retraction}
Let $D_1,\ldots,D_N$ be effective divisors on $X$, and let $Y\to X$
be a log resolution of their union. Include all exceptional divisors
in its SNC divisor, with components $E_j$. For any positive multiple
$v$ of a prime-divisor order over $Y$, set $u_j=v(E_j)$. Then
\begin{equation}\label{eq:retraction}
 v(D_i)=\sum_j u_j\ord_{E_j}(D_i),\qquad
 \A(v)\geq\sum_j u_j\A(E_j).
\end{equation}
If some $u_j$ is nonzero, the monomial order with these weights at the
stratum containing the generic center of $v$ has the same orders on
the $D_i$, no larger discrepancy, and a closed center on $X$
containing that of $v$.
\end{lemma}
\begin{proof}
Each pulled-back divisor is locally a monomial times a unit, which
gives the first equality. The reduced SNC divisor on $Y$ is log
canonical: equivalently, a logarithmic volume form has at most a
simple pole along every prime on a higher smooth model. Thus the
log discrepancy over $Y$ is at least $\sum_j u_j$. Adding the order of
$K_{Y/X}$ gives the second inequality and the claimed discrepancy
of the monomial order. The stratum closure contains the generic
center, so its proper image contains the original closed center.
\end{proof}

\subsection{Thresholds in families}
For a coherent ideal $\mathfrak a$ on a smooth variety, its local log
canonical threshold at $x$ is computed on a log resolution by
\begin{equation}\label{eq:lct-formula}
 \lct_x(\mathfrak a)=
 \min_{x\in\pi(F),\,\ord_F(\mathfrak a)>0}
       \frac{\A(F)}{\ord_F(\mathfrak a)}.
\end{equation}
The minimum is $+\infty$ for the unit ideal. For the zero ideal the
threshold is zero. The same formula follows by change of variables
from local integrability of
$(\sum |g_i|^2)^{-c}$ for generators $g_i$ of $\mathfrak a$.

\begin{lemma}[Closed threshold loci]\label{lem:lct-family}
Let $V$ be a smooth complex projective variety, let $S$ be a complex
algebraic variety, and let $\mathfrak a$ be an ideal on $V\times S$
locally given by algebraic families of generators. For a real number
$c$, the locus
\[
 \{(x,s):\lct_x(\mathfrak a_s)\leq c\}
\]
is Zariski closed. Here $\mathfrak a_s$ denotes the ideal generated
by the restricted generators, with no flatness requirement on its
zero scheme.
\end{lemma}
\begin{proof}
Locally embed the parameter space in a smooth space and extend the
finitely many generators holomorphically to a coordinate neighborhood
of that ambient space. The analytic semicontinuity of complex
singularity exponents applies to
$\varphi=\tfrac12\log\sum |g_i|^2$ there and hence on $S$;
$e^{\varphi}$ is locally H\"older continuous. In local coordinates,
translation of the variable handles a moving marked point.
Consequently the displayed locus is closed in the classical
topology; see \cite[Lemma~3.2]{dk2001}.

It is also algebraically constructible. The following generic-resolution
stratification argument parallels the proof of
\cite[Theorem~3.1]{dk2001}. Resolve the ideal over the
generic point of each irreducible component of $S$ and spread the
resolution over a dense open subset. After shrinking, generic
smoothness makes the fibers and all relevant crossings smooth, so
the fiber maps are log resolutions with fixed orders and
discrepancy coefficients. Formula~\eqref{eq:lct-formula} expresses
the threshold condition as a finite union of images of the
relevant divisors on this piece. These images are algebraic because
the resolution is proper. A finite cover used to separate
components preserves constructibility of the image. Induction on
the closed complement gives a finite stratification and proves
constructibility on $V\times S$. The zero-ideal case is separated
as an algebraic condition. In dimension zero the assertion follows
directly from the vanishing of the generators.

A constructible subset of a complex algebraic variety which is
classically closed is Zariski closed: every Zariski open subset of
an irreducible closed subvariety is classically dense there. This
proves the assertion.
\end{proof}

\subsection{Vanishing and descent}
The cohomological input is Kawamata--Viehweg vanishing
\cite{kawamata1982,viehweg1982}, in the following rounding formulation
\cite[Theorem~0.1]{fujinoKV}.

\begin{theorem}[Kawamata--Viehweg vanishing]\label{thm:kv}
Let $T$ be smooth and projective over $\C$. Let $\Delta$ be an
effective rational SNC divisor with coefficients in $[0,1)$, and
let $D$ be an integral divisor. If $D-(K_T+\Delta)$ is nef and big,
then $H^i(T,\cO_T(D))=0$ for every $i>0$.
\end{theorem}

This is the rounding formulation with
$Q=D-K_T-\Delta$: its fractional support is SNC and
$K_T+\ceil{Q}=D$. In particular, for an integral divisor $M$ and an
effective rational SNC divisor $B$, nefness and bigness of $M-B$
give vanishing for $K_T+M-\floor{B}$, by taking $\Delta=\{B\}$.

We will repeatedly use the following elementary descent observation.
\begin{lemma}\label{lem:descent}
Let $\pi:T\to V$ be a proper birational morphism of smooth varieties,
let $P_V$ be a line bundle on $V$, and let $C$ be an integral divisor
on $T$ whose positive part is exceptional. A section of
$\pi^*P_V+C$ defines a regular section of $P_V$ on $V$.
If a prime $F$ on $T$ maps to a point $x$, the coefficient of $C$
on $F$ is zero, and the section is nonzero generically on $F$, then
the descended section is nonzero at $x$.
\end{lemma}
\begin{proof}
View the section as a rational section of $P_V$ on the common
function field. Its orders on all nonexceptional primes are
nonnegative, since the coefficients of $C$ there are nonpositive.
It is therefore regular in codimension one on the normal variety
$V$, and hence everywhere. Under the additional assumptions its
pullback has order zero along $F$. A local function vanishing at
$x$ has positive order along every prime centered at $x$, so the
descended section cannot vanish at $x$.
\end{proof}
